\chapter{Mengentheoretische Folgerungen}
\label{b06guide:ch:images}

In diesem Kapitel sei stets $f\colon A\rightarrowtail B$.
Weiter seien $U,W\subseteq A$, sofern nicht anders angegeben.

\section{Elemente und Teilmengen zurückgewinnen}

\begin{theorem}[Bildzugehörigkeit und einzelne Fasern]
\label{b06guide:thm:membership}
Für $x\in A$ gilt
\[
 \begin{aligned}
 f(x)\in f[U]&\Longleftrightarrow x\in U,\\
 f^{-1}[\{f(x)\}]&=\{x\}.
 \end{aligned}
\]
Für jedes $b\in B$ ist die Faser leer, falls $b\notin f[A]$,
und sonst eine Einermenge.
\end{theorem}

\begin{theorem}[Gleichheit und Inklusion anhand der Bilder]
\label{b06guide:thm:reflection}
Es gelten
\[
 \begin{aligned}
 f[U]=f[W]&\Longleftrightarrow U=W,\\
 f[U]\subseteq f[W]&\Longleftrightarrow U\subseteq W,\\
 f[U]\subsetneq f[W]&\Longleftrightarrow U\subsetneq W.
 \end{aligned}
\]
Hier bedeutet $\subsetneq$ echte Teilmengeninklusion.
\end{theorem}

\begin{theorem}[Eindeutige Rekonstruktion einer Teilmenge]
\label{b06guide:thm:uniquesubset}
Für jede Teilmenge $V\subseteq f[A]$ gibt es genau eine
Teilmenge $U\subseteq A$ mit $f[U]=V$. Sie lautet
\[
 U=f^{-1}[V].
\]
Die Eindeutigkeit ist eine Folgerung der Injektivität.
\end{theorem}

\section{Mengenoperationen und Durchschnitte}

\begin{theorem}[Exakte Bildregeln]
\label{b06guide:thm:operations}
Es gelten
\[
 \begin{aligned}
 f[U\cap W]&=f[U]\cap f[W],\\
 f[U\setminus W]&=f[U]\setminus f[W],\\
 f[U\mathbin{\triangle}W]&=f[U]\mathbin{\triangle}f[W],\\
 f[A\setminus U]&=f[A]\setminus f[U].
 \end{aligned}
\]
Das rechte Komplement wird relativ zu $f[A]$ gebildet.
Für die Zielmenge $B$ gilt
\[
 B\setminus f[U]
   =(B\setminus f[A])\cup f[A\setminus U].
\]
\end{theorem}

\begin{theorem}[Erkennen von Disjunktheit]
\label{b06guide:thm:disjoint}
\[
 U\cap W=\varnothing
 \quad\Longleftrightarrow\quad
 f[U]\cap f[W]=\varnothing.
\]
Für jede gegebene Familie $(U_i)_{i\in I}$ von Teilmengen von $A$
sind die $U_i$ genau dann paarweise disjunkt,
wenn die $f[U_i]$ paarweise disjunkt sind.
\end{theorem}

\begin{definition}[Durchschnitt relativ zu einer Grundmenge]
Für $E_i\subseteq X$, $i\in I$, sei
\[
 \bigcap_{i\in I}^{X}E_i
    =\{x\in X\mid\forall i\in I\,(x\in E_i)\}.
\]
Der leere Durchschnitt ist in dieser Notation $X$.
\end{definition}

\begin{theorem}[Bilder beliebiger Durchschnitte]
\label{b06guide:thm:intersections}
\label{b06guide:thm:relativeintersection}
Für jede gegebene Familie $(U_i)_{i\in I}$ von Teilmengen von $A$ gilt
\[
 f\left[\bigcap_{i\in I}^{A}U_i\right]
   =\bigcap_{i\in I}^{f[A]}f[U_i].
\]
Für $I\ne\varnothing$ ist dies die gewöhnliche Schnittregel
\[
 f\left[\bigcap_{i\in I}U_i\right]
    =\bigcap_{i\in I}f[U_i].
\]
Die relative Form schließt $I=\varnothing$ ein:
Beide Seiten sind dann $f[A]$.
\end{theorem}